\chapter{Files, other programs and Live Link}
\label{ch:files}

\section{The project}
The Project window is \code{Assets}: nested folders, drag to move, \key{F2} to rename, \key{Delete} to the Recycle Bin
or Trash; references to a moved or renamed material, texture or clip keep working. Scenes are text files
(\code{.scene}); \key{Ctrl+S} saves, \ui{New} or \ui{Open} asks ``Save changes?'' first; the Console's \code{save},
\code{open} and \code{new} do the same.

\section{Import}
\menu{File > Import Model} (or drop files on the window): OBJ + MTL, FBX, glTF, STL, PLY and USD text, with their
materials and textures and, for rigged files, their skeletons, skins, bind poses, shape keys and animations
(\code{Assets/Animations/<model>/}); MikuMikuDance \code{.pmx} / \code{.pmd} characters rigged, toon-shaded and
dancing to \code{.vmd}; VRM avatars with their humanoid names, expressions, spring bones and MToon look; BVH motion
capture; Collada; \code{.blend} through your installed Blender. The \textbf{Import Preview} shows the model first with
its shading switchable (As Imported, PBR, or a toon style) and an Outline switch; Import writes it, Cancel nothing.
The texture finder reads the folders around the model and pairs each bare material with the texture set its name is
likeliest to own.

\section{Export}
\menu{File > Export}: every mesh format Blender 5.2 exports (OBJ, FBX binary or ASCII, glTF \code{.glb} or
\code{.gltf}, STL, PLY, USD) with Blender's options (Selected Only, Apply Modifiers, Scale, Up Axis, Triangulate, Keep
Transforms). glTF and FBX carry a rig, its skin, its clip and the shape keys (visemes, blinks and expressions reach
Unity, VRChat and Unreal); textures go with the file (procedural, painted and generated maps as PNGs, embedded in a
\code{.glb}); \ui{Lightmap UVs} adds a second UV set; \ui{Colliders} and \ui{LOD Levels} are named for Unity, Unreal or
Godot; \ui{Bake Textures} turns procedural materials into images; an export asked for while a remesh or a bind still
works on what it writes waits for it. Every format was checked both ways against Blender 5.2.

\section{Unity projects, VRChat, Godot, Blender, Unreal}
\textbf{Unity:} \menu{File > Import Unity Package\ldots} unpacks a \code{.unitypackage} into
\code{Assets/Imported/<package>/} with Unity's GUIDs kept; a prefab or a \code{.unity} scene comes in with its hierarchy,
meshes, materials (converted by shader: Standard, URP, lilToon, Poiyomi, MToon, UTS2), lights, cameras, colliders,
nested prefabs and overrides, and every other component as an Engine Component you can edit under Unity's names;
\menu{File > Open (Unity Scene)} opens a \code{.unity} as the scene with the look Unity gives it (chapter~\ref{ch:light});
Unity \code{.anim} clips import and export; \menu{File > Export Unity Package\ldots} writes the work back under the same
GUIDs. Nothing of Unity's is ever written over: an edit of a Unity clip or material saves to a Blendity copy beside it.

\textbf{VRChat:} the \textbf{Avatar} window (after the CATS plugin \cite{CATS2026}) has Fix Model, Create Visemes,
Model Options, Pose and Shape tools, Decimation, Bone Root and Bone Merging, Eye Tracking, the Texture Atlas, Merge
Armatures, Attach Mesh, Translate and Export for VRChat, with a checklist against VRChat's requirements
\cite{VRChatCreators2026}; the Profiler's \ui{Budget} ranks a scene or avatar against the PC and Quest performance
ranks; an avatar's VRChat components (the Avatar Descriptor, PhysBones, colliders, contacts, constraints) come in as
Blendity's and go back through the package or Live Link; a world's Scene Descriptor and Light Volumes too.

\textbf{Godot:} \code{.tscn} scenes come in with their nodes, primitives, materials, lights, cameras and collision
shapes and go back with the objects' transforms, scripts kept \cite{GodotDocs2026}. \textbf{Blender:} \code{.blend}
files come in through the installed Blender (\code{blender path} names it); the Blender bridge's \ui{Send to Blendity}
brings a scene over. \textbf{Unreal:} the bridge sends a level through its FBX exporter with each actor's place and
Blueprint class.

\section{Live Link}
Live Link sends the scene to Blender, Unity, Godot or Unreal \emph{as you work}: objects, transforms, meshes (a drag in
Edit Mode streams while it happens), materials, textures by their pixels, lights, cameras, a rig's poses and skins,
clips baked as the engine's own, colliders and LODs, the world and the render camera's look, VRChat's components where
the SDK is there. It starts with Blendity (Preferences $>$ Live Link) and writes where it listens, so a bridge connects
by itself. To put a bridge in: \menu{Window > Live Link > Install a Bridge} $>$ Blender (found and enabled for you), or
Unity / Godot / Unreal Project\ldots\ and pick the project's folder (in Unity: Package Manager $>$ Add package from
disk $>$ the bridge's \code{package.json}). The window lists who is connected. It is one way: edits on the other side
are not sent back. Work on a copy of a project you care about.

\begin{fromnote}{Where it comes from}
The formats follow Blender's importers and exporters \cite{Blender2026} and their specifications
\cite{Khronos2021Gltf, FBXBlender2026, USD2026}; Unity's files follow its manual \cite{UnityManual2026,
UnityPrefabsandtheYAMLsceneformat, UnityExportingapackage}; VRChat's limits and components follow its creator
documentation \cite{VRChatCreators2026}; the Avatar window follows CATS \cite{CATS2026}; Godot's format its
documentation \cite{GodotDocs2026}. Learn: \emph{Assets, Scenes and Files}, \emph{Coordinate Systems and Handedness}.
\end{fromnote}
